\chapter{Thesis Outline}


In the rest of the thesis, we present our work on the Dy-K system.  This work is organized in three additional Parts.

Part \hyperref[pt:questDyKFR]{II} is dedicated to the characterization of the interspecies interaction spectrum of the Dy-K mixture.
\begin{itemize}
	\item In Chapter \ref{Chpt:expIntro},  we first detail the experimental sequence routinely employed in the experiment to produce the doubly degenerate quantum mixture, and the  we further contextualize the experimental findings of this thesis by briefly summarizing the results previously obtained in the experiment
	
	\item Chapter \ref{Chpt:Lowfieldpaper} is dedicated to the results of Ref.~\cite{Ye2022ool}. In this publication we explored the low magnetic field region, below \SI{10}{G}, identifying and characterizing $5$ Feshbach resonances
	
	\item Chapter \ref{Chpt:DipoleandHydro}
	
	\item Chapter \ref{Chpt:AddDyKRes}
\end{itemize}






study of Dy-K system, focus on low magnetic field, individuation of a promising Feshbach resonances, characterization of their properties and in particular of one with a relevant universal range


and characterization of the interaction there.

Part \hyperref[pt:MakingProbing]{III}

\begin{itemize}
	\item Chapter \ref{Chpt:FeshmolIntro}
	
	\item Chapter \ref{Chpt:MolTrap}
	
	\item Chapter \ref{Chpt:MolLoss}
	
\end{itemize}



Part  \hyperref[pt:End]{IV}
